% ===========================================================================
\section{Supplementary Note S5: quality assurance and the audited nitrogen chain}
\label{sec:s5}
% ===========================================================================

The greenhouse-gas and air-pollutant layers of this work draw on the same excreted
nitrogen. That sharing is the system's main structural safeguard, but it also means a
single misplaced fraction propagates into both layers at once. This note records the audit
of the manure nitrogen chain against the 2019 Refinement, the two implementation errors it
found, and the regression tests that now pin the corrections. It is included because the
corrected ledger described in the main text's methods is the outcome of that audit rather
than an inherited assumption, and a reader should be able to see the evidence.

\subsection*{Defect 1: one volatilisation fraction used for two boundaries}

The IPCC separates two volatilisation fractions that an earlier implementation of this
work had conflated. $\mathrm{Frac_{GasMS}}$ (Table~10.22) is the nitrogen lost as
\NHthree{} and \NOx{} \emph{inside} the manure-management system --- housing and storage,
before manure reaches a field --- and it is what reduces the manure available for
application. $\mathrm{Frac_{GASM}}$ (Table~11.3) is the nitrogen lost \emph{after} organic
manure is applied to soil or deposited at grazing; under Tier~1 it drives indirect
\NtwoO{} but does \emph{not} reduce the organic nitrogen input before direct soil
\NtwoO{}.

The earlier implementation used $\mathrm{Frac_{GASM}}$ at the pre-field boundary, where
$\mathrm{Frac_{GasMS}}$ belongs, and then applied it a second time to the surviving
nitrogen. One post-application fraction therefore stood in for a pre-field loss and was
charged twice to the same nitrogen. The corrected chain applies the full Eq.~10.34a term
at the management boundary with per-system Table~10.22 factors, and
$\mathrm{Frac_{GASM}}$ exactly once, post-application. The effect on the national
greenhouse-gas total is $-0.62$~Mt~\COe{}, almost all of it in managed soils, because the
organic nitrogen reaching cropland ($F_{ON}$) falls from \SI{2.721}{} to
\SI{2.335}{}~Mt~N once the pre-field loss is charged at the correct boundary.

\subsection*{Defect 2: an ammonia share applied to a factor already in \NHthree{}-N}

The EMEP/EEA synthetic-fertiliser product factors are published as \emph{ammonia
nitrogen}. An earlier version multiplied the product-weighted value by a further
\NHthree{} share of the volatilised pool, which understates ammonia and simultaneously
leaves the IPCC indirect-\NtwoO{} pool short of its \NOx{} term, so the split stops
conserving nitrogen. The three quantities are now defined and computed separately, as set
out in the main text's air-pollutant methods: $f_{\NHthree{},SF}=0.1608$ drives the ammonia
estimate directly, $f_{NO,SF}=0.007$ is a separate soil process, and only the IPCC
indirect pathway uses their sum, $\mathrm{Frac_{GASF}}=0.1678$, because Eq.~11.9 is
defined on the combined volatilised pool. The net effect on the national total is
$+0.13$~Mt~\COe{}.

\subsection*{Diagnostic configurations}

Because the two chains disagreed in aggregate before they disagreed visibly in any single
stage, the audit compared them stage by stage on a common \NHthree{}-N basis rather than
on totals alone, and the resulting stage table is reported in the main text. That comparison showed the
two chains' national livestock figures differing by only \num{1.24}$\times$ while
individual stages differed by factors of roughly two and three in opposite directions ---
a cancellation that a totals-only check cannot see.

\subsection*{Regression tests}

The corrections are pinned rather than described. The test suite fixes the audited
national total, the per-stage nitrogen balance on both bases, the identity
$\mathrm{Frac_{LossMS}} = \mathrm{Frac_{GasMS}} + \mathrm{Frac_{leachMS}} +
\mathrm{Frac_{N_2MS}} + EF_3$ per species and system, and the requirement that the
manure-management scenarios in the sensitivity and structural-uncertainty tables be
produced by the same full ledger rerun --- a cross-module check added after re-slicing a
single category produced two different national totals for one scenario in two tables of
this paper. Every table in this manuscript is additionally regenerated from the engine and
compared against the committed source, so a number cannot drift from the code that
produced it without a test failing.
